% © 2026 Ing. David Jaroš — CC BY-NC-ND 4.0
%
% This work is licensed under a Creative Commons Attribution-NonCommercial-NoDerivatives
% 4.0 International License (CC BY-NC-ND 4.0).
%
% License History: Earlier drafts (up to v0.3) were released under CC BY 4.0.
% From v0.4 onward, all material is released under CC BY-NC-ND 4.0 to protect
% the integrity of the theoretical work during ongoing academic development.
%
% See LICENSE.md for full license text.

\ifdefined\INCLUDEMODE
\else
  \documentclass[12pt,a4paper]{article}
  \usepackage[a4paper, margin=2.5cm]{geometry}
  \usepackage{amsmath,amssymb,amsthm}
  \usepackage{booktabs}
  \usepackage{tcolorbox}
  \newtheorem{lemma}{Lemma}[section]
  \newtheorem{proposition}[lemma]{Proposition}
  \newtcolorbox{statusbox}[1][]{colback=gray!8!white,colframe=black!60,
    arc=3pt,boxrule=1pt,fonttitle=\bfseries,title=Status Summary,#1}
  \title{Prime Selection Criterion for $p=137$}
  \author{Ing.\ David Jaro\v{s}}
  \date{2026-05-10}
  \begin{document}
  \maketitle
\fi

\section{Arithmetic part: Kronecker criterion}

For odd prime $p$:
\begin{align}
\nu_2(p)
&= 1 + \left(\frac{-4}{p}\right)
 = 1 + \left(\frac{-1}{p}\right)\left(\frac{4}{p}\right)
 = 1 + (-1)^{(p-1)/2}.
\end{align}
Hence
\begin{equation}
\nu_2(p)=2 \iff p\equiv 1\pmod 4,\qquad
\nu_2(p)=0 \iff p\equiv 3\pmod 4.
\end{equation}
This step is purely arithmetic [L0].

\begin{lemma}[Congruence filter on genus-11 prime cluster]
Among primes with $g(X_0(p))=11$, namely $p\in\{131,137,139\}$:
\begin{itemize}
  \item $131\equiv 3\pmod 4 \Rightarrow \nu_2=0$,
  \item $137\equiv 1\pmod 4 \Rightarrow \nu_2=2$,
  \item $139\equiv 3\pmod 4 \Rightarrow \nu_2=0$.
\end{itemize}
Therefore the condition $g(X_0(p))=11$ together with $\nu_2(p)>0$
selects $p=137$ uniquely in this cluster.
\end{lemma}

\section{UBT $U(1)_{EM}$ lift: verification and failure point}

\textbf{Step 1 (required source: \texttt{canonical/interactions/qed.tex}).}
The canonical QED file defines gauge transformations and running
(\texttt{qed.tex}, Sec. ``Gauge Symmetry'' and ``Renormalization''),
but it does not define an explicit map
$C: \Theta\mapsto \Theta^*$ on winding sectors, nor a fixed-point condition
for compactification modes.

\textbf{Step 2.} The modular equivalence
``fixed points of Atkin--Lehner $w_p$ on $X_0(p)$ $\Leftrightarrow \nu_2(p)>0$''
is standard arithmetic geometry [L0], but the bridge
``UBT charge conjugation $C$ $\leftrightarrow$ $w_p$'' is not derived from
$S[\Theta]$ in the cited canonical QED file.

\textbf{Failure location.}
The derivation fails at the physical identification
\begin{equation}
C\text{-fixed winding sector} \Rightarrow w_p\text{-fixed points on }X_0(p)
\end{equation}
because this map is not yet proved in canonical UBT files.
Classification of this bridge: [OPEN].

\begin{proposition}[Conditional physical lift]
If a first-principles UBT derivation establishes
$C\leftrightarrow w_p$ and requires $C$-fixed sectors,
then $\nu_2(p)>0$ follows and therefore $p\equiv 1\pmod 4$.
Under that extra assumption, $p=137$ is selected over $p=139$.
\end{proposition}

\begin{statusbox}
Argument: charge-conjugation selects $p \equiv 1 \pmod 4$\\
Výsledek: OPEN (arithmetic part proved, physical bridge unresolved)\\
$p=137$ je jediné prvočíslo s genus=11 a $p\equiv1\pmod4$ v rozsahu [50,500]: ANO\\
Klasifikace: arithmetic filter [L0], UBT physical lift [OPEN]\\
Dopad: selekce $p=137$ je algebraicky nucena pouze po doplnění otevřeného mostu $C\leftrightarrow w_p$
\end{statusbox}

\ifdefined\INCLUDEMODE
\else
  \end{document}
\fi
